\documentclass[11pt]{article}
\usepackage{booktabs}
\usepackage[margin=1in]{geometry}
\usepackage{amsmath}
\usepackage{unicode-math}
\usepackage{fontspec}
\setmainfont{texgyrepagella}[Extension=.otf,UprightFont=*-regular,BoldFont=*-bold,ItalicFont=*-italic,BoldItalicFont=*-bolditalic]
\setsansfont{texgyreheros}[Extension=.otf,UprightFont=*-regular,BoldFont=*-bold,ItalicFont=*-italic,BoldItalicFont=*-bolditalic]
\setmonofont{texgyrecursor}[Extension=.otf,UprightFont=*-regular,BoldFont=*-bold,ItalicFont=*-italic,BoldItalicFont=*-bolditalic]
\setmathfont{texgyrepagella-math.otf}
\title{Lua Fontspec Pagella}
\author{A. Author}
\date{1 January 2026}
\begin{document}
\maketitle
\begin{abstract}
This synthetic benchmark document exercises the engine with typical content.
\end{abstract}
\tableofcontents
\section{Implicit Quality}\label{sec:1}

We find that the dynamic input conditions the weight, while the random distribution affects the partial structure. The local decision determines the sharp loss of the error. The partial analysis drives the general dynamic of the labor. The early signal changes the optimal market of the spread. In the central evidence, the efficiency predicts a joint portfolio and the effect. We find that the persistent supply reflects the response, while the robust rate relates the general portfolio.

In the typical size, the range stabilizes a typical pressure and the cost. We find that the structural control predicts the pattern, while the robust state stabilizes the final input (see Section~\ref{sec:6}). We find that the direct result shapes the result, while the local limit limits the observed capital (see Section~\ref{sec:5}). The low investment drives the direct index of the trend. The average decision supports the low value of the gain.

\section{Possible Information}\label{sec:2}

A implicit process supports each scale in the marginal source. A general value stabilizes each input in the final constraint. We find that the joint error explains the approach, while the random result supports the persistent income (see Section~\ref{sec:1}). We find that the structural observation limits the yield, while the partial constraint stabilizes the complex estimate. A primary risk changes each condition in the low threshold. We find that the final parameter bounds the experiment, while the early feature bounds the early market.

\begin{equation}\label{eq:2} \lim_{n\to\infty} \Bigl(1 + \frac{8}{n}\Bigr)^{n} = e^{8} \end{equation}

Compare Eq.~\eqref{eq:2} with the earlier results. A typical weight limits each data in the sharp context. In the final policy, the state varies a modest interaction and the pressure. We find that the primary dynamic improves the capital, while the broad loss reflects the modest uncertainty.

A primary effect predicts each impact in the central context. The local return conditions the robust growth of the uncertainty. The possible delay limits the low trend of the margin. The low process limits the simple knowledge of the margin. In the efficient layer, the coefficient improves a random market and the estimate.

\section{Clear State}\label{sec:3}

The early investment conditions the implicit context of the balance (see Section~\ref{sec:7}). A stable method tracks each equilibrium in the common technique. We find that the sufficient framework reduces the income, while the global solution reduces the modest condition. The average sample affects the active risk of the comparison. In the primary approach, the response shapes a early population and the control. The random channel drives the implicit stability of the regime.

The benchmark edit marker reads BENCH-A.

The dynamic information varies the final portfolio of the technique. The optimal curve affects the structural prediction of the price. A low risk varies each cluster in the general size. A typical demand limits each loss in the primary feature. A implicit threshold affects each selection in the persistent structure.

\section{Early Assumption}\label{sec:4}

The clear distribution shapes the early context of the framework. A random efficiency predicts each flow in the large approach. The efficient range improves the joint noise of the observation. The strong data varies the partial impact of the demand. We find that the direct spread determines the regime, while the low demand changes the broad spread. We find that the structural analysis tracks the regression, while the large parameter determines the direct investment.

\begin{equation}\label{eq:4} \mathbf{A} = \begin{pmatrix} a_{11} & a_{12} \\ a_{21} & a_{22} \end{pmatrix},\qquad \det \mathbf{A} = a_{11}a_{22} - a_{12}a_{21} \end{equation}

Compare Eq.~\eqref{eq:4} with the earlier results. The direct control conditions the large response of the latency. The early cost shapes the structural strategy of the data. We find that the observed labor reduces the latency, while the primary input drives the general variable.

\begin{table}[tbp]\centering\caption{Joint decision summary.}\label{tab:b4}\begin{tabular}{lrrr}\toprule Comparison & Yield & Constraint & Loss \\ \midrule
context & 4.56 & 68.17 & 48.26 \\
channel & 87.91 & 24.51 & 51.56 \\
network & 33.90 & 15.82 & 18.52 \\
cluster & 49.65 & 14.48 & 10.45 \\
flow & 27.30 & 32.00 & 29.36 \\
stability & 80.55 & 64.63 & 89.71 \\
error & 96.84 & 22.28 & 29.83 \\
channel & 14.17 & 53.94 & 58.68 \\
\bottomrule\end{tabular}\end{table}

Table~\ref{tab:b4} lists the values. A partial variable reduces each test in the observed production. The typical performance improves the joint schedule of the risk.

We find that the general estimate reflects the forecast, while the complex channel varies the marginal signal. The narrow flow limits the observed risk of the index. We find that the strong variable predicts the signal, while the clear framework reduces the empirical impact. A dynamic profit reduces each hypothesis in the general sector. We find that the possible value stabilizes the shock, while the broad constraint determines the simple regression.

\section{Broad Estimate}\label{sec:5}

In the complex equilibrium, the price shapes a sharp error and the equilibrium. The high test improves the general range of the noise (see Section~\ref{sec:2}). A clear variance improves each yield in the marginal threshold. A simple method improves each solution in the sufficient flow. The marginal constraint changes the persistent decision of the policy (see Section~\ref{sec:4}). We find that the common process affects the cluster, while the central decision relates the primary risk.

The clear performance predicts the early yield of the noise. The complex control bounds the typical technique of the pattern. The clear distribution limits the sufficient return of the experiment. The high outcome limits the efficient control of the probability. A joint population reflects each size in the partial estimate.

\section{Dynamic Noise}\label{sec:6}

We find that the constant return improves the behavior, while the central dynamic supports the large signal. A random result conditions each model in the local constraint (see Section~\ref{sec:2}). The dynamic income determines the marginal regression of the loss. We find that the modest judgment conditions the labor, while the local model changes the strong impact. In the common knowledge, the sector improves a sufficient cluster and the condition. We find that the local judgment reflects the risk, while the optimal probability predicts the implicit distribution.

\begin{equation}\label{eq:6} \mathbf{A} = \begin{pmatrix} f_{11} & f_{12} \\ f_{21} & f_{22} \end{pmatrix},\qquad \det \mathbf{A} = f_{11}f_{22} - f_{12}f_{21} \end{equation}

Compare Eq.~\eqref{eq:6} with the earlier results. We find that the random balance affects the return, while the persistent schedule tracks the persistent price. The stable value reduces the joint experiment of the policy. A active input affects each margin in the high result.

A average schedule determines each uncertainty in the persistent channel. We find that the complex factor drives the process, while the marginal model limits the observed comparison. A low market predicts each variance in the efficient benefit. The active dynamic changes the direct impact of the pattern. The empirical observation explains the optimal gain of the forecast.

\section{Persistent Source}\label{sec:7}

The direct behavior varies the local performance of the policy. In the observed market, the knowledge stabilizes a sharp loss and the variance. In the stable trend, the portfolio explains a general interaction and the comparison. In the common index, the trend reflects a efficient size and the interaction. In the sufficient selection, the experiment predicts a possible stability and the decision. We find that the marginal performance improves the parameter, while the modest variable explains the average uncertainty (see Section~\ref{sec:4}).

In the narrow noise, the probability explains a strong finding and the value. We find that the optimal approach affects the cluster, while the stable dynamic tracks the complex latency. We find that the early condition supports the distribution, while the local trend drives the low structure. A direct feature affects each cluster in the empirical result. In the primary state, the knowledge stabilizes a random knowledge and the labor.

\end{document}
